\section{Conclusiones}
\label{sec:conclusiones}

\paragraph{Elección del filtro.} De los cuatro diseños que cumplen la plantilla, se implementó el pasabajos Chebyshev~I de orden~3. Es una función
todo polos, de orden impar y $Q$ moderado, que se realiza con una red RC con seguidor y un Sallen--Key de ganancia unitaria en un solo TL072, y tiene
casi \qty{18}{\decibel} de margen a \qty{4}{\kilo\hertz}. Su desventaja, el retardo de grupo variable cerca del borde de paso, se trató en el
Nivel~III.

\paragraph{El margen de diseño fue la decisión clave.} La función de la plantilla tiene exactamente \qty{1}{\decibel} de rizado, y con componentes
reales solo cumpliría alrededor del \qty{\vRendSinMargen}{\percent} de los filtros construidos. Sintetizar en cambio una función con
\qty{\vApD}{\decibel} de rizado y borde en \qty{\vFpD}{\hertz} cambia la atenuación a \qty{4}{\kilo\hertz}, que sobraba, por tolerancia en la banda
de paso, que faltaba. El prototipo lo confirma: sus capacitores se apartan hasta un \qty{10}{\percent} del nominal y aun así su atenuación máxima en
la banda de paso es de \qr{2}{\mApMaxMed}{\decibel}. Con los mismos desvíos, el filtro sin margen tendría \qr{2}{\mApSinMargen}{\decibel} y no
cumpliría la plantilla.

\paragraph{El prototipo cumple la plantilla y el modelo lo explica.} La atenuación a \qty{4}{\kilo\hertz} es de \qr{1}{\mAcuatrokMed}{\decibel},
con \qty{11}{\decibel} de margen sobre los \qty{24}{\decibel} exigidos. El modelo ajustado a la ganancia y la fase medidas describe la medición
dentro de la incertidumbre ($\chi^2_\nu = \nr{2}{\mChiCuadRed}$). Los capacitores estimados a partir de ese ajuste, junto con las resistencias
medidas y las sensibilidades de la \cref{sec:sensibilidad}, predicen una variación de la ganancia a \qty{1}{\kilo\hertz} de
\qr{2}{\mDGunokPred}{\decibel}; la medida es de \qr{2}{\mDGunokMed}{\decibel}. Como anticipaba el análisis de sensibilidad, los capacitores, y no
las resistencias ni el AO, determinan cuánto se aparta el filtro real del nominal.

\paragraph{Cadena de validación.} Cada etapa se contrastó con la anterior y todas coincidieron dentro de su incertidumbre: la derivación simbólica
con SymPy, el circuito ideal en LTspice (diferencia menor que \qty{1e-4}{\decibel}), el modelo del TL072 (cuyo único efecto apreciable en la
banda de interés es el GBW finito) y la medición. Las discrepancias que aparecieron tuvieron causas identificables y documentadas: la compresión de
formas de onda de LTspice en la respuesta al escalón, la fuga del Sallen--Key de ganancia unitaria por encima de \qty{40}{\kilo\hertz} y el piso
de medición del osciloscopio en la banda de rechazo lejana.

\paragraph{Nivel~II.} El Monte Carlo de \nNRun{} corridas distintas por escenario en LTspice estima un rendimiento del \qty{\nRendLtDos}{\percent}
con los componentes usados, en acuerdo con Python (\qty{\nRendPyDos}{\percent}) y con la predicción por sensibilidades
(\qty{\nRendLinDos}{\percent}). Con tolerancias de hasta \qty{10}{\percent}, el filtro se comporta de forma prácticamente lineal y la fórmula de
propagación de sensibilidades basta para estimar la dispersión, aunque no la forma de la distribución. Todas las fallas se deben al rizado, por
el desvío de los capacitores: con capacitores de \qty{\pm 5}{\percent} cumple el \qty{100}{\percent} de los circuitos, mientras que mejorar las
resistencias aporta poco.

\paragraph{Nivel~III.} Un pasatodo de 2.º orden con un solo AO, diseñado para el prototipo y no para el filtro nominal, redujo la variación del
retardo de grupo en la banda de paso de \qr{4}{\eMJFil}{\milli\second} a \qr{4}{\eMJCas}{\milli\second}: una mejora de \num{\eMMejora} veces,
con un intervalo del \qty{90}{\percent} de \numrange{\eMMejoraInf}{\eMMejoraSup}. La métrica sin modelo lo confirma: el desvío máximo de la fase
respecto de la fase lineal bajó de \qr{1}{\eMFilDevMax}{\degree} a \qr{1}{\eMCasDevMax}{\degree}. En la respuesta a la onda cuadrada, la
deformación pasa de concentrarse después del flanco a repartirse a ambos lados, como corresponde a un sistema de fase casi lineal. El precio es un
retardo constante adicional de medio milisegundo y una ganancia $k = \num{\eMApK}$; el error de pasatodo remanente es un valle de
\qr{1}{\eMApValle}{\decibel} en $f_0$, con el que la cascada sigue dentro de la plantilla en la banda de paso. La lección principal del nivel es práctica: el óptimo exige $f_0$ y $Q$ con errores de décimas de por ciento, y los
capacitores solo se conocían al \qty{4}{\percent}. El ecualizador alcanzó la mejora porque se sintonizó midiendo: el preset fija la magnitud sin
mover $f_0$ ni $Q$, y la resistencia menor de cada par corrige $f_0$ y $Q$ a partir de los capacitores efectivos deducidos de la medición. La
simulación con los capacitores leídos con el multímetro da una mejora mucho menor, lo que muestra que esa sintonía no se puede reemplazar con
cálculo.

\paragraph{Limitaciones y trabajo futuro.}
\begin{itemize}
  \item \textbf{Capacitores.} Sin medidor LCR, sus valores reales se estimaron en forma indirecta. Un medidor LCR permitiría verificar esas
        estimaciones y calcular el ecualizador sin iteraciones. Para una producción, capacitores de \qty{\pm 5}{\percent} eliminarían prácticamente las
        fallas sin cambiar la topología.
  \item \textbf{Banda de rechazo lejana.} Por encima de unos \qty{4.5}{\kilo\hertz} la medición queda limitada por el piso del instrumento. Medirla
        requiere más amplitud de entrada, promediado o un analizador con más rango dinámico.
  \item \textbf{Gran señal.} El ensayo de recorte no pudo hacerse con la amplitud máxima del generador; puede completarse bajando la alimentación.
  \item \textbf{Ganancia del ecualizador.} La cascada tiene \qr{1}{\eLTEcGananciaDb}{\decibel} de ganancia. Si hiciera falta ganancia unitaria, se
        puede agregar una etapa no inversora o redistribuir la ganancia en el filtro.
  \item \textbf{Respuesta a la onda cuadrada.} La captura final no se tomó en modo de alta resolución, lo que limita la incertidumbre del
        sobrepico a uno o dos puntos porcentuales.
\end{itemize}

En síntesis, el filtro construido cumple la plantilla con margen, y cada resultado medido se explica con el diseño, la sensibilidad y la
simulación. Los dos niveles optativos agregan la cuantificación estadística del rendimiento y una mejora medida de más de diez veces en la
linealidad de fase. Todo el trabajo es reproducible desde los notebooks, los netlists y los datos del repositorio.
